\documentclass[letterpaper,journal]{IEEEtran}
\usepackage{amsmath,amsfonts}
\usepackage{algorithmic}
\usepackage{algorithm}
\usepackage{array}
\usepackage[font=small, labelfont=rm,textfont=rm]{caption}
\usepackage[font=small, labelfont=rm,textfont=rm]{subcaption}
\usepackage{textcomp}
\usepackage{stfloats}
\usepackage{url}
\usepackage{verbatim}
\usepackage{graphicx}
\usepackage{cite}
\usepackage{multirow}
\usepackage{array}
\usepackage{booktabs}
\usepackage[table]{xcolor}
\usepackage{makecell}
\usepackage[T1]{fontenc}
\usepackage[bitstream-charter]{mathdesign}
\renewcommand{\rmdefault}{ptm}
\usepackage{times}
% \hyphenation{op-tical net-works semi-conduc-tor IEEE-Xplore} % removed template sample
% updated with editorial comments 8/9/2021

\newcommand{\sigmafunc}{\mathrm{sigm}} 
\begin{document}


\title{TerraSAP: Spatially-Aware Prompt-based Framework for Few-Shot Class-Incremental Learning in Remote Sensing Image Classification}


\author{
    Jiaxing Zeng,
    Yifeng Tan,~\IEEEmembership{Student Member,~IEEE,}
    Lina Yang,~\IEEEmembership{Member,~IEEE,}
    Siwei Zhang,
    and Lianhui Liang,~\IEEEmembership{Member,~IEEE}
        % <-this % stops a space



\thanks{Jiaxing Zeng is with the School of Computer, Electronics and Information, Guangxi University, Nanning 530004, China, and is also with the Institute of Systems Engineering and the Department of Engineering Science, Faculty of Innovation Engineering, Macau University of Science and Technology, Macau 999078, China (email: zasakraea@gmail.com).}

\thanks{Lina Yang is with the School of Computer, Electronics and Information, Guangxi University, Nanning 530004, China (email: lnyang@gxu.edu.cn).}

\thanks{Yifeng Tan and Lianhui Liang are with the School of Electrical Engineering, Guangxi University, Nanning 530004, China (email: tyf1282331092@163.com; lianglh0308@126.com).}

\thanks{Siwei Zhang is with the Institute of Systems Engineering and the Department of Engineering Science, Faculty of Innovation Engineering, Macau University of Science and Technology, Macau 999078, China (email: swzhang@must.edu.mo).}

\thanks{Corresponding authors: Lina Yang and Siwei Zhang}
}


% The paper headers
\markboth{}%
{Shell \MakeLowercase{\textit{et al.}}: A Sample Article Using IEEEtran.cls for IEEE Journals}

\maketitle


% --- Abstract ---
\begin{abstract}
The core challenge of Few-Shot Class-Incremental Learning (FSCIL) lies in continuously learning new categories with extremely few samples without forgetting old knowledge. However, the performance of general FSCIL solutions is often limited when directly applied to the remote sensing domain. This limitation stems from two unique challenges of remote sensing images, namely complex and explicit spatial structural relationships and significant intra-class heterogeneity. Most existing methods fail to adequately model these characteristics. To address this issue, this paper proposes an efficient framework designed specifically for remote sensing scenarios, called TerraSAP. To tackle the aforementioned challenges, the framework systematically addresses the unique problems of remote sensing FSCIL through two core designs. First, to overcome the limitation of lacking spatial priors, we propose a Spatially-Aware Prompt Module. It explicitly injects the inherent spatial topological structure priors of remote sensing images into deep features through deterministic position encoding and attention mechanisms. Second, to overcome the learning difficulties caused by complex visual characteristics, we design a synergistic mechanism for new-class learning and knowledge retention. This mechanism includes two core parts. First, it enhances the learning efficiency for sparse new classes through an NCAC-based prototype initialization strategy. Second, it achieves a fine balance between retaining old knowledge and learning new knowledge with an innovative Dual-Path Exponential Moving Average (EMA) mechanism. Experimental results show that TerraSAP achieves an average accuracy of 84.52\% on the NWPU-RESISC45 dataset, with a significant improvement in new-class learning ability. It also achieves efficient parameter utilization with a 47.10\% reduction in total parameters, a 93.20\% reduction in trainable parameters, and a 27.80\% increase in inference speed.
\end{abstract}

% --- Keywords --- 
\begin{IEEEkeywords}
Few-shot class-incremental learning, remote sensing image classification, catastrophic forgetting, prompt learning, parameter-efficient fine-tuning.
\end{IEEEkeywords}


% --- Main Text ---
\section{INTRODUCTION}

\IEEEPARstart{R}{emote} sensing images (RSI), as a core data source for Earth observation, play an indispensable role in urban planning, environmental monitoring, disaster assessment, and military reconnaissance~\cite{cheng2017, ma2019, zhu2017grsm}. Deep Neural Networks (DNNs), with their powerful feature representation capabilities, have become a cornerstone technology for RSI interpretation~\cite{cheng2017, ma2019, zhu2017grsm}. However, the great success of these methods is mostly built upon a static and idealized assumption that training data for all classes are complete and sufficient before model deployment. This static setting starkly contrasts with the dynamic nature of real-world remote sensing applications. In the real world, remote sensing data are continuously accumulated over time, and new land-cover classes, such as new types of buildings, newly discovered natural landforms, emerge dynamically as observation tasks evolve and expand. At the same time, the samples for these new classes are often extremely scarce. This is mainly caused by a combination of factors, including the limited accessibility of observation platforms, short imaging windows, and high annotation costs. This practical dilemma poses an urgent demand for the model's Lifelong Learning ability, i.e., the ability to continuously learn new knowledge without forgetting old knowledge.

To address the above issues, Class-Incremental Learning (CIL) \cite{rebuffi2017, deLange2022, parisi2019, vandeven2019} has emerged. CIL aims to enable models to continuously learn from a stream of data with ever-emerging new categories. The core challenge lies in overcoming Catastrophic Forgetting~\cite{kirkpatrick2017}---the phenomenon where a model rapidly forgets previously learned knowledge while learning new information.

However, traditional CIL methods typically assume that samples for new classes are abundant, which contradicts the reality of remote sensing applications. This more stringent real-world setting has given rise to Few-Shot Class-Incremental Learning (FSCIL) \cite{tao2020}. FSCIL integrates the challenges of both CIL and Few-Shot Learning (FSL) \cite{snell2017, finn2017, wang2020}. It must not only combat catastrophic forgetting but also avoid overfitting on extremely sparse samples, presenting a significant challenge.

To tackle the challenges of FSCIL, the academic community has conducted extensive exploration. In the general image domain, mainstream methods mitigate forgetting through techniques such as replay~\cite{lopez2017, aljundi2019}, regularization~\cite{kirkpatrick2017, zenke2017}, and parameter isolation~\cite{rusu2016, xu2018}. They also adapt to new categories by constructing continuously evolving classifiers~\cite{zhang2021} or reserving feature space for new classes~\cite{zhou2022}. The Attention-based Self-Adaptive Prompt (ASP)~\cite{liu2024} method, by introducing attention-aware adaptive prompts, encourages task-invariant prompts to capture shared knowledge while leveraging an information bottleneck to transfer knowledge from old classes to new ones, achieving significant results in mitigating catastrophic forgetting. In addition, meta-learning based approaches, for example, MAML-style initialization, have been explored to improve rapid adaptation under few-shot constraints~\cite{finn2017, ravi2017, hospedales2021meta}, and knowledge distillation strategies are often combined with exemplar replay to stabilize old-class decision boundaries during incremental updates~\cite{rebuffi2017, li2017lwf, hinton2015distill, cheraghian2021skd}.

Despite the considerable progress of these methods, they still face two key limitations when confronted with the complex characteristics of remote sensing images:

\begin{enumerate}
    \item \textbf{Inability to integrate remote sensing-specific spatial prior knowledge}. Land-cover objects in remote sensing images, for example buildings, roads, and farmland, are not isolated pixel clusters but follow specific geospatial distribution patterns; for example, harbors are always accompanied by water and bridges inevitably span rivers. This spatial context information is a crucial clue for distinguishing between classes with highly similar visual features \cite{li2022, huang2022}. However, existing general-purpose FSCIL methods largely lack explicit modeling of such spatial structure priors, relying solely on appearance features, which can easily lead to confusion in complex scenes.
    \item \textbf{Large intra-class variations and subtle inter-class similarities}. Remote sensing objects of the same class can exhibit vastly different appearances (intra-class variation). This phenomenon is the result of a combination of factors such as imaging conditions (illumination, season, sensor angle), geographical environment variability, and resolution differences. Conversely, objects of different classes may possess highly overlapping spectral or texture features due to similar materials or structures (inter-class similarity) \cite{clemente2017, brusch2011}. This characteristic imposes extremely demanding requirements on the model's representation ability and the precision of its decision boundaries, a task that traditional methods struggle with in few-shot scenarios.
\end{enumerate}

\textcolor{blue}{To address these issues more pointedly, researchers have also begun to explore new methods specifically for remote sensing images. For example, the Cosine Prototype Learning (CPL) framework~\cite{zhao2023} leverages cosine prototypes to cope with the significant intra-class variation observed in SAR imagery. The HDCPAA model~\cite{li2025} optimizes the classification boundaries for optical remote sensing images using hyperdimensional computing to handle subtle inter-class similarities. Furthermore, the RSSC method~\cite{nwpu_cankao} mitigates representation drift in remote sensing scene FSCIL through gradient-guided multi-scale feature collaboration, while PCOFA~\cite{zeng2022pcofa} enhances few-shot incremental performance in aerial scenes through feature adaptation and continuous prototype optimization. There have also been explorations in different modalities, such as hyperspectral class-incremental learning with distillation and SAR incremental recognition methods~\cite{xu2022rs,chen2019,zhang2025pfs3f,liang2025dbmlla,zhang2025msgfus}. Meanwhile, SAM-based change detection under limited supervision~\cite{gao2025samchange} and efficient open-vocabulary segmentation tailored to remote sensing scenes~\cite{cao2025openvocab} illustrate complementary directions that exploit foundation models and broad vocabularies for Earth observation. Although these methods have achieved some success, they often focus on solving one of the challenges or rely on single-modality/weak spatial relationship assumptions, failing to fully leverage the rich spatial context information embedded in remote sensing images to address both challenges simultaneously. Consequently, their performance in complex scenes still has room for improvement. To this end, we propose TerraSAP (Terra Spatially-Aware Prompt-based FSCIL), an efficient, modular, and synergistic learning framework designed specifically for remote sensing scenarios. We build upon the advanced ASP method~\cite{liu2024} but fundamentally reconstruct its core prompt generation and knowledge update mechanisms. The goal is to deeply integrate remote sensing-specific spatial priors into the learning process and significantly enhance parameter efficiency. Our main contributions are summarized as follows.}

\begin{enumerate}
\item We propose a Spatially-Aware Prompt Module that deterministically encodes absolute and relative 2D positions into prompts, injecting remote sensing spatial priors into a frozen ViT for more discriminative features. It explicitly models spatial topology and is more parameter-efficient than variational designs.

\item We introduce a New Class-Aware Classifier equipped with a single-step prototype initialization strategy that leverages spatial context to initialize novel-class prototypes, improving few-shot learning under high variance and mitigating intra-class heterogeneity in remote sensing.

\item We design a Dual-Path Exponential Moving Average (EMA) strategy that maintains stable and plastic paths to balance stability--plasticity, consolidating old knowledge while adapting to new classes and reducing forgetting under high inter-class similarity.

\item We validate on NWPU-RESISC45, UCM, and MSTAR, showing higher accuracy and cross-modal generalization with fewer parameters and faster inference than state-of-the-art baselines.
\end{enumerate}

\section{PRELIMINARIES}
This section provides the necessary background on prompt-based methods, and prototypical networks. The definitions and formulations are primarily based on the ASP paper \cite{liu2024} to set the stage for our proposed method.

\subsection{Prompt-based Methods in Vision Tasks}
Prompt-based methods in vision tasks typically use a pre-trained Vision Transformer (ViT) as the backbone $f_\theta$ and freeze its parameters $\theta$ during training to preserve the generalization capabilities acquired from pre-training. A ViT model consists of multiple multi-head self-attention (MSA) layers. Let the input to the $l$-th MSA layer be $h^l\in\mathbb{R}^{L_h\times D}$. The output of this layer is:
%
\begin{equation}
\begin{aligned}
\mathrm{MSA}(h^l) &= \operatorname{Concat}\big(h_1^l,\ldots,h_m^l\big)\,W^O,\\
h_i^l &= \operatorname{Attention}\!\big(h^lW_i^Q,\,h^lW_i^K,\,h^lW_i^V\big)\\
\operatorname{Attention} &=\operatorname{Softmax}\!\left(\frac{(h^lW_i^Q)(h^lW_i^K)^{\top}}{\sqrt{d_k}}\right)(h^lW_i^V),
\end{aligned}
\end{equation} where $W^O, W_i^Q, W_i^K, W_i^V$ are projection matrices. The terms $h^l W_i^Q$, $h^l W_i^K$, and $h^l W_i^V$ are the head-specific query, key, and value representations derived from $h^l$; $m$ is the number of attention heads; and $d_k$ denotes the dimensionality of each key vector.

In prompt-based methods, Prompt Tuning introduces a small number of trainable parameters $p^l\in\mathbb{R}^{L_p\times D}$ as prompts at layer $l$, which are prepended to the input $h^l$:
\begin{equation}
f_{\mathrm{ProT}}(p^l,h^l)=\mathrm{MSA}([p^l;\;h^l]),
\end{equation} where $[\,\cdot\;;\;\cdot\,]$ denotes concatenation along the sequence length dimension. Before entering the first layer of the ViT, an input image is first divided into patches and converted into a sequence-like representation $x^e\in\mathbb{R}^{L_x\times D}$. For image classification tasks, a class token $cls\in\mathbb{R}^{1\times D}$ is concatenated along with the visual prompts to form the input to the ViT module:
\begin{equation}
x^p=[cls;\;p^0;\;x^e]
\end{equation}

\subsection{Prototypical Network}
Prototypical Networks~\cite{snell2017} are a common method in few-shot learning. They compute a mean feature vector $c_k$ for each class $k$, which serves as the class prototype:
\begin{equation}
c_k=\frac{1}{N_k}\sum_{y_i=k} f(x_i),
\end{equation} where $N_k$ is the number of samples in class $k$, and the output of the class token can be used as the image embedding. For a classification task with $K$ classes, $W=[c_0,\ldots,c_K]$ is used as the linear classifier. An input sample is classified based on the softmax probability over its similarity to the prototypes:
\begin{equation}
P(y=k\mid x)\propto c_k^\top f_{\theta,p}(x)
\end{equation}

Following previous work, prototypes for new classes are appended to $W$ during incremental learning to cover all seen classes.

\section{PROPOSED METHOD}

\subsection{Problem Definition and Setting}
We follow the standard training pipeline for FSCIL. We divide the data into $T+1$ sessions, $\{\mathcal{D}_0,\mathcal{D}_1,\dots,\mathcal{D}_T\}$, with corresponding class sets $\{\mathcal{C}_0,\mathcal{C}_1,\dots,\mathcal{C}_T\}$, where the class sets are mutually exclusive, i.e., $\mathcal{C}_i \cap \mathcal{C}_j = \emptyset, \forall i \neq j$.

The training process is divided into two phases. The first is the \textbf{base session} ($t=0$), where the model is trained on the base class dataset $\mathcal{D}_0$, which contains ample samples. This is followed by a series of \textbf{incremental sessions} ($t=1, \dots, T$), where in each session $t$, the model is exposed to only a few ($K$-shot) samples of new classes from $\mathcal{D}_t$. After each session, the model's performance is evaluated on the test set of all seen classes $\bigcup_{i=0}^{t}\mathcal{C}_i$.

\begin{figure*}[!t]
    \centering
    \includegraphics[width=\textwidth]{TerraSAP_Framework_overview.png}
    \caption{\textbf{TerraSAP framework overview.} 
    \textbf{(Left) Base session.} A frozen ViT receives an RSI image while the Spatially-Aware Prompt Module (SAPM) injects spatial priors (absolute 2D positional encoding and relative position bias) to produce task-specific prompts. The trainable parts are SAPM and the cosine classifier. At the end of the base session, two EMA memories are initialized for Dual-Path EMA (DPEMA): the stability path $\mathrm{AvgTSP}_{\text{orig}}$ (from non-spatially-modulated prompts) and the plasticity path $\mathrm{AvgTSP}_{\text{mod}}$ (from spatially-modulated prompts).
    \textbf{(right) Incremental session.} For each novel class, NCAC performs single-step prototype initialization using support features and spatial context $\{f_i\},\{\mathrm{ctx}_i\}$: compute $\bar{p}_c, s_c, \gamma_c$ and obtain the modulated prototype $\hat{p}_c$, which is appended to the classifier. Meanwhile, DPEMA updates $\mathrm{AvgTSP}_{\text{orig}}$ and $\mathrm{AvgTSP}_{\text{mod}}$ via EMA and fuses them with weight $\alpha_{\mathrm{ema}}$ to balance stability and plasticity across sessions. Frozen vs.\ trainable modules and data flow are indicated in the figure.}
    \label{fig:framework_overview}
\end{figure*}

\subsection{Framework Overview}
Our TerraSAP framework is built upon a pre-trained Vision Transformer (ViT) \cite{dosovitskiy2021, khan2022transformers}, with its backbone parameters kept frozen throughout the learning process to achieve extreme parameter efficiency. For ease of exposition, we will use TIP, Template-Initialized Prompt, and TSP, Task-Specific Prompt~\cite{liu2024}. The framework is technically based on the advanced ASP method \cite{liu2024}, but we have designed three synergistic core innovative modules to address the unique challenges of remote sensing images, as illustrated in Fig.~\ref{fig:framework_overview}. The overall procedure is summarized in Algorithm~\ref{alg:terrasap_fscil}:

\subsubsection{Spatially-Aware Prompt Module (SAPM)} To address the rich spatial structure information in remote sensing images, this module generates and injects prompts carrying spatial topology priors, enabling the frozen ViT backbone to perceive and utilize the spatial relationships between ground objects, thereby obtaining more discriminative feature representations.

\subsubsection{New Class-Aware Classifier (NCAC)} To achieve rapid and efficient learning when new-class samples are extremely scarce, this module adopts an \emph{NCAC-based prototype initialization}: it first computes a mean prototype from the support features, and then modulates it using spatial context to obtain the final prototype. This two-step process provides a better initial decision boundary in a single initialization step.
\subsubsection{Dual-Path Exponential Moving Average (DPEMA) Synergistic Mechanism} To address the problem of catastrophic forgetting (stability), this mechanism achieves a finer balance between retaining old knowledge and learning new knowledge by maintaining in parallel an original knowledge base biased towards stability and a modulated knowledge base biased towards plasticity.

These three modules work in synergy. SAPM first enhances the quality of features. Based on these high-quality features, NCAC and DPEMA jointly solve the stability-plasticity dilemma in incremental learning and ultimately achieve efficient FSCIL in remote sensing scenarios.

\begin{algorithm}[!t]
\caption{TerraSAP FSCIL Training Pipeline}
\label{alg:terrasap_fscil}
\begin{algorithmic}[1]
\REQUIRE Sessions $\{\mathcal{D}_t\}_{t=0}^{T}$, frozen ViT $f_{\theta}$, hyperparameters $(\lambda_{\mathrm{anchor}}, \gamma_0, \ldots)$
\ENSURE Trained parameters $\Theta$, EMA states $\{A^{(t)}_{\mathrm{orig}}, A^{(t)}_{\mathrm{mod}}\}_{t=0}^{T}$
\STATE \textbf{Base Session ($t=0$)}
\STATE Initialize SAPM, NCAC, DPEMA; classifier weights $\mathbf{W}^{(0)}$; optimizer $\mathcal{O}$
\STATE $P_{\mathrm{spatial}} \gets \mathrm{SAPM\_Generate}$ per Eq.~\eqref{eq:attention_scores}
\STATE $P_{\mathrm{mod}} \gets \mathrm{ThreeBranchMod}(P_{\mathrm{spatial}})$ per Eq.~\eqref{eq:three_branch_modulation}
\REPEAT
    \STATE Sample $(\mathbf{x},\mathbf{y}) \sim \mathcal{D}_0$
    \STATE $\mathbf{F} \gets \mathrm{ViT\_Forward}(\mathbf{x}, P_{\mathrm{mod}})$
    \STATE $L \gets L_{\mathrm{CE}}(\mathbf{F},\mathbf{y}) + \lambda_{\mathrm{anchor}} L_{\mathrm{anchor}}(\mathbf{F},\mathbf{y})$ per Eqs.~\eqref{eq:base_loss_actual}, \eqref{eq:anchor_loss}
    \STATE $\Theta \gets \mathcal{O}\text{-Update}(\Theta, \nabla_{\Theta} L)$
\UNTIL{validation convergence}
\STATE $A_{\mathrm{orig}}^{(0)} \gets \mathrm{TSP}^{(0)}$, $A_{\mathrm{mod}}^{(0)} \gets \mathrm{TSP}_{\mathrm{mod}}^{(0)}$
\STATE \textbf{Incremental Sessions ($t=1,\dots,T$)}
\FOR{$t = 1$ to $T$}
    \STATE Choose support subset $\mathcal{D}_t^{\mathrm{support}} \subset \mathcal{D}_t$ and new-class set $\mathcal{C}_t$
    \STATE $\text{TSP}^{(t)} \gets \mathrm{Mean}(\{\text{TSP}_i\}_{(\mathbf{x}_i,\mathbf{y}_i)\in\mathcal{D}_t^{\mathrm{support}}})$ using frozen SAPM
    \FOR{$c \in \mathcal{C}_t$}
        \STATE $(\{\mathbf{f}_i\}, \{\mathrm{ctx}_i\}) \gets \mathrm{ExtractFeatures}(\mathcal{D}_{t,c}^{\mathrm{support}}, f_{\theta})$
        \STATE $\bar{\mathbf{p}}_c \gets \frac{1}{|\mathcal{D}_{t,c}^{\mathrm{support}}|}\sum_i \mathbf{f}_i$ per Eq.~\eqref{eq:base_prototype}
        \STATE $\overline{\mathrm{ctx}}_c \gets \mathrm{Mean}(\{\mathrm{ctx}_i\})$, $\mathbf{s}_c \gets \tanh(\Psi(\overline{\mathrm{ctx}}_c))$ per Eq.~\eqref{eq:spatial_weight}
        \STATE $\gamma_c \gets \gamma_0\,\sigmafunc(\operatorname{std}(\overline{\mathrm{ctx}}_c))$ per Eq.~\eqref{eq:modulation_strength}
        \STATE $\hat{\mathbf{p}}_c \gets \mathrm{AdaptiveMod}(\bar{\mathbf{p}}_c, \mathbf{s}_c, \gamma_c)$ per Eq.~\eqref{eq:adaptive_modulation}
    \ENDFOR
    \STATE $\mathbf{W}^{(t)} \gets [\,\mathbf{W}^{(t-1)};\{\hat{\mathbf{p}}_c\}_{c\in\mathcal{C}_t}\,]$
    \STATE $(A_{\mathrm{orig}}^{(t)}, A_{\mathrm{mod}}^{(t)}) \gets \mathrm{EMA\_Update}(A_{\mathrm{orig}}^{(t-1)}, A_{\mathrm{mod}}^{(t-1)}, \text{TSP}^{(t)})$ per Eq.~\eqref{eq:ema_update}
    \STATE $\operatorname{AvgTSP}^{(t)} \gets \mathrm{DualEMA\_Fuse}(A_{\mathrm{orig}}^{(t)}, A_{\mathrm{mod}}^{(t)})$ per Eq.~\eqref{eq:dual_ema_fusion}
    \STATE Evaluate on seen classes $\bigcup_{i=0}^{t} \mathcal{C}_i$
\ENDFOR
\end{algorithmic}
\end{algorithm}


\subsection{Spatially-Aware Prompt Module}

We apply a Manhattan-distance bucketization between token pairs, prompt--prompt and prompt--patch, to derive relative position indices, which modulate attention logits additively via a learned bias. This is implemented inside the prompt encoder and does \emph{not} alter the frozen ViT backbone.
General prompt learning methods often treat prompts as an unordered set of vectors, which neglects the crucial spatial structure priors in remote sensing images. To address this issue, we designed SAPM. Its core idea is to enable the prompts themselves to encode and transmit spatial topology information, thereby guiding the ViT backbone to focus on the spatial context between ground objects during feature extraction. 

\textcolor{blue}{
Unlike methods like ASP~\cite{liu2024} that use a variational encoder, we adopt a deterministic encoder architecture to directly generate spatially-aware prompts. This design differs fundamentally from standard position-aware attention mechanisms in RS Transformers, which require training backbone parameters. Instead, SAPM injects spatial priors into a frozen ViT via prompt tokens. This deterministic design avoids the stochastic sampling noise of variational encoders, ensuring robust topological representation, e.g., for roads and bridges, in few-shot regimes where stability is paramount. This design choice is based on two key considerations. }

First, the spatial relationships in remote sensing images are deterministic in nature. The spatial topology of ground objects, such as roads and airports, follows physical and geographical laws and is therefore highly deterministic. In contrast, the random sampling perturbations introduced by a variational encoder might corrupt these fine-grained spatial structure priors. Second, a deterministic encoder is more stable in few-shot scenarios, avoiding the sampling uncertainty introduced by inaccurate fitting of the complex posterior distribution in variational inference, thus reducing the risk of overfitting.

We construct spatially-aware prompts by fusing absolute and relative position information. First, we assign a unique 2D sinusoidal positional encoding to each prompt token, giving it an absolute sense of geometric coordinates. Then, in the standard self-attention computation, we introduce a relative position bias term $S_{\text{rel}}$ in Eq. (\ref{eq:attention_scores}). This bias term is a learnable embedding that adjusts based on the relative displacement between prompt tokens, thereby explicitly modeling the prior knowledge of proximity and distance. The attention scores are computed as:
{\setlength{\abovedisplayskip}{6pt}\setlength{\belowdisplayskip}{6pt}%
\begin{align}
\label{eq:attention_scores}
\mathbf{S}
\,=\, &\; \frac{\mathbf{Q}\mathbf{K}^\top}{\sqrt{d_k}} \, +\, \omega_{\mathrm{rel}}\,\mathbf{S}_{\mathrm{rel}},\\
&\; [\mathbf{S}_{\mathrm{rel}}]_{ij} \,=\, \Phi\!\left(\Delta_{ij}\right),\qquad \Delta_{ij}=\lVert \mathbf{p}_i-\mathbf{p}_j\rVert_1,
\end{align}} where $\mathbf{Q}$ and $\mathbf{K}$ stack the head-specific query and key matrices $h^lW_i^Q$ and $h^lW_i^K$ introduced in the MSA formulation, $d_k$ is the attention head key dimension, $\mathbf{S}_{\mathrm{rel}}\in\mathbb{R}^{L\times L}$ is the relative bias, $\omega_{\mathrm{rel}}\in[0,1]$ is a scalar weight for spatial relations, $\Delta_{ij}$ is the Manhattan distance on the 2D grid, and $\Phi: \mathbb{R}_{\ge 0}\to\mathbb{R}$ is a learnable mapping in Eq.~(\ref{eq:attention_scores}).

Finally, a lightweight MLP network applies a non-linear transformation to the attention-computed prompts, further enhancing their ability to model complex spatial relationships.

To enhance the robustness of the generated prompts while maintaining parameter efficiency, we introduce a lightweight three-branch modulation strategy, see Eq. (\ref{eq:three_branch_modulation}). This mechanism is designed to stabilize the prompt features by combining three complementary operations, a spatial smoothing branch ($\operatorname{AvgPool1D}$) to reduce high-frequency noise, a semantic stabilization branch, $\operatorname{LayerNorm}$ with mild additive noise, to regularize the feature distribution, and a consistency branch that preserves the original prompt structure. These branches are dynamically combined using trainable scales, controlled by an overall modulation strength. In each Transformer layer of the ViT, we concatenate the generated spatially-aware prompts with the image patch tokens. The spatial information is propagated to all tokens through the self-attention mechanism, after which the prompt tokens are trimmed off. This plug-and-play approach efficiently injects spatial information without altering the ViT backbone structure. Their formulas are as follows:
\begin{align}
\label{eq:three_branch_modulation}
P_{\mathrm{sp}} &= \operatorname{AvgPool1D}(P,\,k{=}3),\\
P_{\mathrm{se}} &= \operatorname{LN}(P) + \epsilon,\quad \epsilon\sim\mathcal{N}(0, 0.01^2),\\
P_{\mathrm{mix}} &= \frac{w_s\,P_{\mathrm{sp}} + w_m\,P_{\mathrm{se}} + w_c\,P}{w_s + w_m + w_c},\\
P_{\mathrm{mod}} &= (1-\eta)\,P + \eta\,P_{\mathrm{mix}},
\end{align} where $w_s, w_m, w_c > 0$ are trainable scales, and $\eta$ controls the modulation strength. In practice, we enforce non-negativity of $w_s,w_m,w_c$ via a smooth constraint.
\subsection{Intelligent Adaptation and Knowledge Retention for Incremental Tasks}
After obtaining more discriminative spatially-enhanced features, we address how to use these features to efficiently learn new classes while preventing catastrophic forgetting of old classes. We tackle this through the synergistic design of NCAC and DPEMA.

Building upon the spatially enhanced features produced by the prompt module, the core challenge shifts to leveraging this rich information for robust new class learning. To this end, we directly incorporate the spatial context ($\mathrm{ctx}_c$) generated by the prompts into classifier initialization, as detailed below.
Standard classifier extension methods, such as random initialization or simple prototype averaging, are inefficient and unstable in few-shot settings. To achieve rapid and accurate new class learning, we propose NCAC. Its core is a NCAC-based prototype initialization that allows the classifier to find a better initial decision boundary in one step.

Unlike the standard prototype method, which only calculates the mean of the support set sample features, our method incorporates spatial context information extracted from the spatially-aware prompt module when calculating new class prototypes. Specifically, we designed an adaptive modulation strategy, their formula is as follows: 

\begin{align}
\label{eq:base_prototype}
\bar{\mathbf{p}}_c &= \frac{1}{N_c}\sum_{i:\,y_i=c} \mathbf{f}_i,\\
\label{eq:spatial_weight}
\mathbf{s}_c &= \tanh\!\left(\Psi(\mathrm{ctx}_c)\right),\\
\label{eq:modulation_strength}
\gamma_c &= \gamma_0 \, \sigmafunc\!\left(\operatorname{std}(\mathrm{ctx}_c)\right),\\
\label{eq:adaptive_modulation}
\hat{\mathbf{p}}_c &= \lambda\,\big(\bar{\mathbf{p}}_c + \gamma_c\,\mathbf{s}_c\big)
                 + (1-\lambda)\, \bar{\mathbf{p}}_c \odot \big(1 + \gamma_c\, \mathbf{s}_c\big),
\end{align} 
where $N_c$ is the number of support samples in class $c$, $\mathrm{ctx}_c$ denotes the class-related context derived from spatial prompts, $\Psi: \mathbb{R}^{d_{ctx}} \to \mathbb{R}^{d}$ aligns the context to the feature dimension $d$, $\Phi: \mathbb{R}_{\ge 0} \to \mathbb{R}$ maps relative distances to bias values in attention, $\sigmafunc$ denotes the sigmoid function, and $\lambda\in[0,1]$ balances additive/multiplicative paths.
This dual modulation not only allows the prototype to shift its position in the feature space but also reshapes its feature distribution, thereby better accommodating the high intra-class variance commonly present in remote sensing data. Moreover, the modulation strength $\gamma_c$ is designed to vary with the standard deviation of the spatial context. The more diverse the spatial context, the stronger the correction required for the base prototype. Finally, all class weights and image features are L2-normalized, and classification is performed via cosine similarity calculation:
\begin{align}
\label{eq:cosine_classifier}
\hat{\mathbf{F}} \,=\, \frac{\mathbf{F}}{\lVert \mathbf{F}\rVert_2},\quad
\hat{\mathbf{W}} \,=\, \frac{\mathbf{W}}{\lVert \mathbf{W}\rVert_2},\qquad
\mathbf{Z} \,=\, \tau\, \hat{\mathbf{F}}\, \hat{\mathbf{W}}^{\top},
\end{align} where $\mathbf{F}\in\mathbb{R}^{B\times d}$ and $\mathbf{W}\in\mathbb{R}^{C\times d}$, $\tau>0$ is a learnable temperature parameter, and class probabilities are obtained as $\operatorname{softmax}(\mathbf{Z})$.

\textcolor{blue}{Excellent classifier initialization alone is not enough to resist catastrophic forgetting in long-sequence tasks. A knowledge retention mechanism is crucial. We argue that in remote sensing FSCIL tasks, traditional single-path EMA mechanisms (such as the one used in ASP~\cite{liu2024}, Eq. 11), which use a single momentum coefficient to update a single knowledge base, struggle to find an optimal trade-off between adapting to large intra-class variations and retaining old knowledge. Furthermore, unlike generic dual-path CL methods that maintain fast and slow model weights, our DPEMA maintains two distinct Prompt Memories: one for raw visual stability ($\operatorname{AvgTSP}_{\text{orig}}$) and one for spatially-modulated plasticity ($\operatorname{AvgTSP}_{\text{mod}}$). This explicitly decouples the types of knowledge rather than just the speed of updating. To achieve a finer balance, DPEMA decouples the knowledge accumulation process, maintaining two independent historical prompt means in parallel:}

Stability Path, $\operatorname{AvgTSP}_{\text{orig}}$, performs an exponential moving average (EMA) update on the original, non-spatially modulated task prompts, aiming to preserve a more general and stable historical knowledge base.
Plasticity Path, $\operatorname{AvgTSP}_{\text{mod}}$, performs an EMA update on the task prompts after they have been modulated by spatial context, aiming to capture more task-relevant and adaptive dynamic knowledge.

Both paths follow the standard EMA update, applied to their respective streams:
\begin{align}
\label{eq:ema_update}
A_{\text{orig}}^{(t)} &= \beta_{\text{ema}}\, A_{\text{orig}}^{(t-1)} + (1-\beta_{\text{ema}})\, \text{TSP}_{\text{orig}}^{(t)},\\
A_{\text{mod}}^{(t)}  &= \beta_{\text{ema}}\, A_{\text{mod}}^{(t-1)}  + (1-\beta_{\text{ema}})\, \text{TSP}_{\text{mod}}^{(t)},
\end{align} where $\beta_{\text{ema}}\in(0,1)$ is the momentum coefficient; $\text{TSP}_{\text{orig}}^{(t)}$ denotes the mean of the current session's non-modulated task prompts, and $\text{TSP}_{\text{mod}}^{(t)}$ denotes the mean after spatial modulation.

During the model's forward pass, the final global average prompt $\operatorname{AvgTSP}$ used to guide feature extraction is a weighted fusion of the knowledge bases from these two paths:
\begin{equation}
\begin{aligned}
\operatorname{AvgTSP}^{(t)} &= \alpha_{\mathrm{ema}}\, \operatorname{AvgTSP}_{\text{orig}}^{(t)} \\
&\quad + (1-\alpha_{\mathrm{ema}})\, \operatorname{AvgTSP}_{\text{mod}}^{(t)},
\end{aligned}
\label{eq:dual_ema_fusion}
\end{equation} where $A_{\text{orig}}^{(t)}$ and $A_{\text{mod}}^{(t)}$ are the historical means from the two paths in Eqs.~(\ref{eq:ema_update})--(\ref{eq:dual_ema_fusion}). When $\beta_{\text{ema}}\to 1$, the update is conservative and stable; when $\beta_{\text{ema}}\to 0$, the update is rapid and plastic. In our implementation, the fusion parameter $\alpha_{\mathrm{ema}}\in[0,1]$ is a fixed hyperparameter that forms a convex combination of the two memories.

\subsection{Optimization Objective and Training Strategy}
We will discuss the model's training process and optimization objective together.

Our optimization objective follows the core design of ASP \cite{liu2024} and consists of the standard cross-entropy loss $L_{\text{CE}}$ and the anchor constraint loss $L_{\text{anchor}}$. We adopt a deterministic spatial encoder and do not introduce additional variational regularization, thus keeping the objective function concise and stable in the few-shot setting.

In the base session ($t=0$), the trainable modules (SAPM and the cosine classifier) are optimized through multiple iterations, with the loss function:
\begin{align}
L = L_{\text{CE}} + \lambda_{\text{anchor}}\,L_{\text{anchor}} \label{eq:base_loss_actual},
\end{align}where $L_{\text{CE}}$ is the standard cross-entropy loss, and $L_{\text{anchor}}$ is an anchor constraint loss that enhances intra-class compactness by minimizing the cosine distance between the features of samples within a class and their class anchor. For reproducibility, we provide the explicit form consistent with the implementation:
\begin{equation}
\label{eq:anchor_loss}
L_{\text{anchor}} \,=\, \frac{1}{|\mathcal{B}|}\sum_{(\mathbf{x}_i,\,y_i)\in\mathcal{B}} \Big(1 - \cos\big(\mathbf{f}_i,\,\mathbf{a}_{y_i}\big)\Big),
\end{equation}where $\mathbf{f}_i$ is the sample feature, and $\mathbf{a}_c$ is the representative anchor feature for class $c$, which is the feature of the sample within the batch that has the highest cosine similarity to the class center (mean feature). It is calculated by first finding the class center $\bar{\mathbf{f}}_c=\frac{1}{|\mathcal{B}_c|}\sum_{j:y_j=c}\mathbf{f}_j$, and then selecting the closest sample $j^{*}=\arg\max_{j:y_j=c} \cos\big(\mathbf{f}_j,\,\bar{\mathbf{f}}_c\big)$ as the anchor $\mathbf{a}_c=\mathbf{f}_{j^{*}}$. This loss maximizes the intra-class cosine similarity on the unit sphere, which is compatible with (\ref{eq:cosine_classifier}).

In the incremental sessions ($t>0$), to avoid overfitting on a very limited number of samples, we adopt a single-step, non-iterative update strategy. All trained modules are frozen, and the classifier weights for new classes are initialized through a single computation based on the new classes' support set using our proposed new class-aware strategy. This stage involves no iterative optimization.

\section{EXPERIMENTS AND ANALYSIS}
\subsection{Implementation Details}
\subsubsection{Training Configuration And Environment}
All experiments were conducted on a server equipped with an NVIDIA RTX 5090 GPU (32GB VRAM), using PyTorch 2.9.0 and CUDA 12.8.
We use a frozen ViT-B/16~\cite{dosovitskiy2021, liu2024} as the backbone. In the base session, we train for 40 epochs with SGD, using a momentum of 0.9 and a weight decay of $5\times10^{-4}$. The initial learning rate is 0.007 and follows a cosine decay schedule. We set the batch size to 16. The prompts contain six tokens, and TIP is initialized to zero. For data augmentation, we apply random resized cropping to $224\times224$ pixels and random horizontal flipping with probability $0.5$. During evaluation, we resize images to $256$ pixels on the shorter side and then take a center crop to $224\times224$ pixels.

\textcolor{blue}{To ensure a fair and rigorous comparison, we strictly standardized the experimental setting. Unless otherwise specified, e.g., CPL, all comparison methods were reproduced or adapted using the same frozen ViT-B/16 backbone pre-trained on ImageNet-21k. We applied identical data preprocessing pipelines and consistent training protocols for the base session to eliminate performance variances caused by backbone differences. This allows us to attribute performance gains specifically to the incremental learning mechanisms rather than feature extractor capacity.}


\subsubsection{Evaluation Metrics And Comparison Agreements} We use Top-1 accuracy and average accuracy (the arithmetic mean of accuracies across all sessions) as the primary metrics. We also report new class accuracy (performance only on newly added classes) and harmonic accuracy (the harmonic mean of old and new class accuracies). Inference time is measured for a single image with a batch size of 1 and FP32 precision, averaged over 200 runs after 15 warm-up iterations.

\subsection{Dataset Introduction}

\subsubsection{Dataset Description} To comprehensively evaluate the performance and generalization ability of our method, we selected three datasets with distinct characteristics in terms of the number of classes, image modality, and scene complexity. NWPU-RESISC45\cite{cheng2017} is a large-scale optical remote sensing scene classification dataset with 45 classes. UCM\cite{ucm} is another widely used optical remote sensing scene dataset with 21 classes. The MSTAR\cite{mstar} dataset is a Synthetic Aperture Radar (SAR) dataset containing images of 10 classes of ground military targets.

\textcolor{blue}{To directly address Reviewer~4's concern regarding urban-domain diversity, we additionally process three hyperspectral benchmarks that capture complex city layouts: So2Sat-LCZ42 (Sentinel-2 multi-spectral imagery labeled with Local Climate Zones), Pavia University (acquired by ROSIS, 103 spectral bands), and the Houston 2013 IEEE GRSS Data Fusion Challenge scene (CASI sensor, 144 bands). These datasets complement NWPU/UCM by focusing on densely built environments and broaden the modality coverage beyond SAR by introducing high-dimensional optical spectra.}

\subsubsection{Dataset Configuration} We followed the experimental splits from RSSC \cite{nwpu_cankao} for the NWPU-RESISC45 and UCM datasets. The 45-class NWPU-RESISC45 dataset was divided into 25 base classes and four incremental sessions (five classes per session). We divided the UCM dataset into 12 base classes and nine incremental classes that are further subdivided into three sessions (three incremental classes per session, for a total of 21 classes). For the MSTAR dataset, we followed the split from CPL \cite{zhao2023}, with four base classes and six incremental sessions (one class per session). For all incremental sessions, we used a uniform 5-shot setting.

We evaluated the effectiveness of TerraSAP on three remote sensing datasets, namely NWPU-RESISC45\cite{cheng2017}, UCM \cite{ucm}, and MSTAR\cite{mstar}. Fig.~\ref{fig:dataset_samples} shows sample examples from these datasets.

\begin{figure}[!t]
\centering
% Unified triple subfigures for consistent layout
\subfloat[NWPU-RESISC45 dataset samples]{\includegraphics[width=\linewidth]{NWPU_RESISC45_samples.png}}\par\medskip
\subfloat[UCM dataset samples]{\includegraphics[width=\linewidth]{UCMerced_samples.png}}\par\medskip
\subfloat[MSTAR dataset samples]{\includegraphics[width=\linewidth]{MSTAR_samples.png}}
\caption{Sample images from benchmark datasets (a) NWPU-RESISC45, (b) UCM, (c) MSTAR.}
\label{fig:dataset_samples}

\end{figure}

\textcolor{blue}{For the new urban hyperspectral benchmarks we adopt reviewer-aligned splits: So2Sat-LCZ42 uses the Base7 + (5 sessions $\times$ 2 classes) schedule with five support samples per new class; Houston 2013 follows a Base7 + (4 sessions $\times$ 2 classes) design that explicitly separates residential, transportation, and recreational structures; and Pavia University adopts a Base5 + (4 sessions $\times$ 1 class) plan because only nine categories are available. Each cube is compressed to a pseudo-RGB representation via 3-component PCA, we slide a $64\times64$ window centered on annotated pixels, and bilinearly upsample the resulting patches to $224\times224$ to remain compatible with the frozen ViT-B/16 backbone. Detailed preprocessing scripts will be released to facilitate reproducibility.}

\begin{table}[!t]
\centering
\caption{\textcolor{blue}{New urban hyperspectral benchmarks and FSCIL splits.}}
\label{tab:urban_datasets}
\begin{tabular}{@{}lccc@{}}
\toprule
Dataset & Modality (bands) & FSCIL Split & \#Classes \\
\midrule
So2Sat-LCZ42 & Sentinel-2 (13) & Base7 + 5$\times$2 Inc. & 17 \\
Pavia University & ROSIS (103) & Base5 + 4$\times$1 Inc. & 9 \\
Houston 2013 & CASI (144) & Base7 + 4$\times$2 Inc. & 15 \\
\bottomrule
\end{tabular}
\end{table}

\subsection{Performance Comparison and Analysis}

\textcolor{blue}{To provide a clear overview, we group the representative state-of-the-art methods into categories:}

\subsubsection{\textcolor{blue}{Metric \& Feature Adaptation}}

\textcolor{blue}{These methods focus on learning robust distance metrics or translating features to new distributions.}

\textcolor{blue}{C-FSCIL \cite{hersche2022cfscil} is a constraint-based method that optimizes a hyperdimensional embedding space.}

\textcolor{blue}{FeTrIL (WACV 2023)~\cite{petit2023fetril} combines a fixed feature extractor with a pseudo-feature generator to translate base class features into new class representations for classifier training.}

\subsubsection{\textcolor{blue}{Prototype}}
\textcolor{blue}{Prototype methods maintain class centers for classification.}

\textcolor{blue}{CPL \cite{zhao2023} is a cosine prototype learning method tailored for SAR target recognition, leveraging domain-specific properties.}

\textcolor{blue}{ALICE (ECCV 2022)~\cite{peng2022alice} utilizes angular penalty losses and mixup strategies to improve the discriminability of class prototypes and margins.}

\subsubsection{\textcolor{blue}{Distillation \& Regularization}}
\textcolor{blue}{These methods mitigate forgetting through constraints on the feature space or logits.}

\textcolor{blue}{CEC \cite{zhang2021} progressively updates a cosine classifier with a graph-based context and distillation.}

\textcolor{blue}{FACT \cite{zhou2022} leverages forward-compatible training to reserve feature space for future classes.}

\textcolor{blue}{BiDist (CVPR 2023)~\cite{zhao2023bidist} employs bidirectional distillation between the old and new models to maintain stability and plasticity.}

\textcolor{blue}{LIMIT (TPAMI 2023)~\cite{zhou2023limit} generates pseudo-features for old classes to compensate for data imbalance during incremental steps.}

\textcolor{blue}{SAVC \cite{song2023savc} introduces semantic-aware virtual contrastive constraints to regularize representations.}

\textcolor{blue}{RSSC \cite{nwpu_cankao} uses multi-scale collaboration regularization specifically for remote sensing scene classification.}

\subsubsection{\textcolor{blue}{Prompt}}
\textcolor{blue}{Prompt methods adapt the backbone with lightweight learnable tokens.}

\textcolor{blue}{ASP \cite{liu2024} learns attention-aware adaptive prompts with an information bottleneck to transfer knowledge.}

\begin{table*}[!t]
\centering
\caption{\textcolor{blue}{Comparison with state-of-the-art methods on NWPU-RESISC45 (Top-1 Accuracy, \%).}}
\label{tab:nwpu_sota}
\begin{tabular}{@{}llrrrrrr@{}}
\toprule
Method (ID) & Reference & S1 & S2 & S3 & S4 & S5 & \textbf{Average} \\
\midrule
CEC & CVPR 2021 \cite{zhang2021} & 93.50 & 82.00 & 78.40 & 70.00 & 67.10 & 78.20 \\
C-FSCIL & CVPR 2022 \cite{hersche2022cfscil} & 93.65 & 79.20 & 76.50 & 75.45 & 68.20 & 78.60 \\
FACT & CVPR 2022 \cite{zhou2022} & 93.90 & 84.10 & 77.30 & 73.00 & 66.20 & 78.90 \\
ALICE & ECCV 2022 \cite{peng2022alice}& 94.15 & 84.00 & 77.60 & 72.80 & 68.30 & 79.37 \\
FeTrIL & WACV 2023 \cite{petit2023fetril}& 94.25 & 85.20 & 79.00 & 72.50 & 68.20 & 79.83 \\
BiDist & CVPR 2023 \cite{zhao2023bidist}& 94.40 & 82.50 & 78.40 & 76.30 & 69.10 & 80.14 \\
LIMIT & TPAMI 2023 \cite{zhou2023limit} & 94.65 & 85.90 & 79.50 & 71.80 & 70.60 & 80.49 \\
SAVC & CVPR 2023 \cite{song2023savc} & 94.30 & 84.00 & 76.00 & 74.50 & 69.20 & 79.60 \\
RSSC & IEEE TGRS 2024 \cite{nwpu_cankao} & 94.80 & 85.90 & 82.80 & 71.80 & 71.80 & 81.42 \\
\midrule
ASP & ECCV 2024 \cite{liu2024} & \textbf{94.55} & 86.79 & 80.46 & 77.03 & 76.16 & 83.00 \\
\textbf{TerraSAP} & \textbf{Ours} & 93.87 & \textbf{86.80} & \textbf{82.39} & \textbf{81.12} & \textbf{78.43} & \textbf{84.52} \\
\bottomrule
\end{tabular}
\end{table*}

\textcolor{blue}{\textbf{Performance and analysis on the NWPU-RESISC45 dataset.}}
\textcolor{blue}{As shown in Table~\ref{tab:nwpu_sota}, TerraSAP achieves a state-of-the-art average accuracy of 84.52\%. This performance exceeds the strong baselines such as LIMIT~\cite{zhou2023limit} and BiDist~\cite{zhao2023bidist}, which also utilize the powerful ViT backbone. The observed gain is plausibly attributable to deterministic spatially-aware prompts. In contrast to general FSCIL methods like FeTrIL~\cite{petit2023fetril} or ALICE~\cite{peng2022alice} that focus on feature distribution manipulation, our method explicitly injects spatial priors, yielding more stable and discriminative features for complex remote sensing scenes. The effect is most pronounced in the first incremental session (S2), where new-class accuracy improves significantly, indicative of stronger early-session adaptation.}
\textcolor{blue}{Beyond early sessions, we also observe consistent robustness in later sessions, where the average accuracy remains competitive despite the compounding class space. This suggests that spatial priors help maintain well-separated decision boundaries even as new classes accumulate. Qualitatively, categories with complex man-made structures benefit the most, as spatial prompts emphasize layout and connectivity cues. Compared to RSSC, which is also tailored for remote sensing, TerraSAP achieves superior retention, demonstrating that explicitly modeling spatial context is more effective than regularization alone in handling large intra-class variations.}

% --- UCM TABLE ---
\begin{table*}[!t]
\centering
\caption{\textcolor{blue}{Comparison with state-of-the-art methods on UCM (Top-1 Accuracy,\%)}}
\label{tab:ucm_sota}
\begin{tabular}{@{}llrrrrr@{}}
\toprule
Method (ID) & Reference & S1 & S2 & S3 & S4 & \textbf{Average} \\
\midrule
CEC & CVPR 2021 \cite{zhang2021} & 96.90 & 77.90 & 74.80 & 70.80 & 80.10 \\
C-FSCIL & CVPR 2022 \cite{hersche2022cfscil} & 96.50 & 80.10 & 78.30 & 72.10 & 81.75 \\
FACT & CVPR 2022 \cite{zhou2022} & 98.20 & 81.80 & 73.12 & 70.00 & 80.78 \\
ALICE & ECCV 2022 \cite{peng2022alice} & 97.40 & 84.30 & 73.20 & 70.50 & 81.35 \\
FeTrIL & WACV 2023 \cite{petit2023fetril} & 97.60 & 79.30 & 77.10 & 72.92 & 81.73 \\
BiDist & CVPR 2023 \cite{zhao2023bidist} & 97.80 & 81.50 & 78.80 & 70.50 & 82.15 \\
LIMIT & TPAMI 2023 \cite{zhou2023limit} & 97.90 & 82.70 & 78.20 & 71.40 & 82.55 \\
SAVC & CVPR 2023 \cite{song2023savc} & 97.70 & 82.00 & 75.10 & 72.72 & 81.88 \\
RSSC & IEEE TGRS 2024 \cite{nwpu_cankao} & 98.30 & \textbf{87.40} & \textbf{80.30} & \textbf{77.00} & \textbf{85.75} \\
\midrule
ASP & ECCV 2024 \cite{liu2024} & 97.29 & 77.65 & 75.52 & 69.89 & 80.09 \\
\textbf{TerraSAP} & \textbf{Ours} & \textbf{98.33} & 85.10 & 79.62 & 74.61 & 84.41 \\
\bottomrule
\end{tabular}
\end{table*}

\textcolor{blue}{\textbf{Performance and analysis on the UCM dataset.}}
\textcolor{blue}{On the UCM dataset (Table~\ref{tab:ucm_sota}), comparisons become tighter as the ViT backbone already provides strong initial features for this relatively small dataset. TerraSAP attains a competitive average accuracy of 84.41\%, significantly outperforming the ASP baseline and general SOTA methods like LIMIT~\cite{zhou2023limit} and BiDist~\cite{zhao2023bidist}. It performs on par with the strong competitor RSSC. The relative strength of RSSC is consistent with the dataset's more regular scenes and smaller label space, where directly regularizing feature-space geometry is highly effective. By contrast, our prompt-based approach emphasizes adaptive representation shaping.}
\textcolor{blue}{Nevertheless, TerraSAP yields clear gains in S1 and S3, indicating that spatial priors still help stabilize new-class insertion under limited support. Gains concentrate on categories where spatial arrangement is key, for example, runways versus roads. Overall, while general regularization methods work well in structured environments, our spatial prompt-based method provides a versatile option that remains top-tier even when scene variation is modest.}

\begin{table*}[!t]
\centering
\caption{\textcolor{blue}{Comparison with state-the-art methods on MSTAR (Top-1 Accuracy, \%). \textbf{General methods use ViT-B/16; CPL uses SAR-specific optimization.}}}
\label{tab:mstar_sota}
\begin{tabular}{@{}llrrrrrrrr@{}}
\toprule
Method (ID) & Reference & S1 & S2 & S3 & S4 & S5 & S6 & S7 & \textbf{Avg} \\
\midrule
CEC & CVPR 2021 \cite{zhang2021} & 92.50 & 76.80 & 72.60 & 62.10 & 58.80 & 51.10 & 48.52 & 66.06 \\
C-FSCIL & CVPR 2022 \cite{hersche2022cfscil} & 92.90 & 87.50 & 74.40 & 68.50 & 62.90 & 57.20 & 52.69 & 70.87 \\
FACT & CVPR 2022 \cite{zhou2022} & 92.80 & 74.20 & 69.50 & 65.20 & 58.90 & 54.70 & 52.23 & 66.79 \\
ALICE& ECCV 2022 \cite{peng2022alice} & 93.20 & 83.90 & 70.80 & 60.30 & 59.40 & 56.10 & 48.59 & 67.47 \\
FeTrIL & WACV 2023 \cite{petit2023fetril} & 93.50 & 84.10 & 75.80 & 63.20 & 57.00 & 55.20 & 46.92 & 67.96 \\
BiDist & CVPR 2023 \cite{zhao2023bidist} & 94.10 & 88.90 & 77.00 & 70.30 & 64.00 & 61.10 & 58.68 & 73.44 \\
LIMIT & TPAMI 2023 \cite{zhou2023limit} & 93.80 & 80.50 & 71.60 & 65.20 & 60.70 & 56.30 & 52.10 & 68.60 \\
SAVC & CVPR 2023 \cite{song2023savc} & 93.40 & 81.80 & 70.10 & 63.90 & 59.50 & 55.30 & 49.97 & 67.71 \\
RSSC & IEEE TGRS 2024 \cite{nwpu_cankao} & 95.20 & 88.20 & 82.50 & 72.00 & 66.40 & 63.20 & 57.78 & 75.04 \\
CPL & IEEE TGRS 2023 \cite{zhao2023} & \textbf{99.89} & \textbf{90.60} & 76.96 & 72.92 & 70.39 & \textbf{66.71} & 60.36 & \textbf{76.83} \\
\midrule
ASP & ECCV 2024 \cite{liu2024} & 92.55 & 83.94 & 72.70 & 63.83 & 59.00 & 57.34 & 54.08 & 69.06 \\
\textbf{TerraSAP} & \textbf{Ours} & 94.79 & 86.82 & \textbf{80.70} & \textbf{76.32} & \textbf{70.45} & 66.08 & \textbf{61.71} & 76.70 \\
\bottomrule
\end{tabular}
\end{table*}

\begin{figure*}[!t]
\centering
\subfloat[NWPU-RESISC45]{\includegraphics[width=0.32\textwidth]{sota_sessions_nwpu.png}}\hfill
\subfloat[UCM]{\includegraphics[width=0.32\textwidth]{sota_sessions_ucm.png}}\hfill
\subfloat[MSTAR]{\includegraphics[width=0.32\textwidth]{sota_sessions_mstar.png}}
\caption{Session-wise accuracy line charts for the three datasets.}
\label{fig:sota_session_lines}
\end{figure*}

\textcolor{blue}{\textbf{Cross-modal generalization on the MSTAR dataset.}}
\textcolor{blue}{Session-wise performance across incremental sessions is summarized in Fig.~\ref{fig:sota_session_lines}. To assess cross-modal generalization from optical imagery to SAR, we evaluate on the MSTAR dataset (Table~\ref{tab:mstar_sota}). Even when powerful ViT baselines are applied, TerraSAP achieves an average accuracy of 76.70\%, consistently outperforming recent general methods like LIMIT~\cite{zhou2023limit} and FeTrIL~\cite{petit2023fetril}. The prototype-based CPL, designed specifically for SAR imaging physics, attains 76.83\% and serves as a strong domain-specific benchmark.}
\textcolor{blue}{We note that the relative advantages are consistent with method specialization. CPL leverages domain-specific prototype separability, whereas our framework demonstrates competitive performance without any SAR-specific adaptation, bridging the gap solely through spatial-aware prompting. Notably, the largest gains over general methods appear in earlier sessions, e.g., S2, where spatial prompts guide the model to exploit shape and configuration cues that remain relevant after the modality shift.}
\textcolor{blue}{Later-session performance remains stable, suggesting that spatial priors, averaged by DPEMA, retain cross-session consistency without heavy rehearsal. These results indicate that structure-aware prompts capture modality-agnostic spatial patterns, offering a robust solution that outperforms standard ViT-based incremental learning techniques on SAR data.}

\textcolor{blue}{\textbf{Urban hyperspectral benchmarks.}}
\textcolor{blue}{Beyond SAR, we now evaluate TerraSAP on the So2Sat-LCZ42, Pavia University, and Houston 2013 datasets summarized in Table~\ref{tab:urban_datasets}. The exact preprocessing mirrors Reviewer~4's suggestion: we apply PCA to compress the spectral cubes to pseudo-RGB patches, use a $64\times64$ sliding window centered on labeled pixels, and upsample to $224\times224$ before feeding the frozen ViT-B/16 backbone. We keep the same 5-shot incremental setting and reproduce the ASP baseline under identical schedules. Due to space limitations, the full quantitative curves are reported in the supplementary material (Tables~S1--S3), but they follow the same trends observed on NWPU/UCM: TerraSAP consistently improves new-class accuracy in the sessions where elongated man-made structures (roads, highways, railways) appear, while maintaining old-class stability thanks to DPEMA.}
\textcolor{blue}{These additional benchmarks demonstrate that TerraSAP scales to high-dimensional optical spectra and densely built scenes, directly addressing Reviewer~4's request for diversity. They also highlight that deterministic spatial prompts act as a modality-agnostic prior: once the hyperspectral cubes are mapped to pseudo-RGB patches, no further architectural changes are required and the frozen ViT continues to operate in the same regime.}

\subsection{Ablation Study}
We conducted systematic ablation studies on the NWPU-RESISC45, UCM, and MSTAR datasets to quantify the individual and synergistic contributions of each module.
Visual Validation of Spatially-Aware Effectiveness. To visually validate the effectiveness of SAPM, we used both ScoreCAM and Grad-CAM to generate attention heatmaps, comparing the model's focus patterns on remote sensing images with strong spatial structural relationships, with and without SAPM. Fig.~\ref{fig:scorecam_bridge} shows the ScoreCAM visualization for the bridge class, providing key visual evidence for our core argument.

\begin{figure}[!t]
    \centering
    \includegraphics[width=\linewidth]{Bridge_ScoreCAM_blocks10.png}
    \caption{Bridge class visualization using ScoreCAM (target layer: last ViT block attention; ReLU normalization; bilinear upsampling) with and without spatial awareness \cite{wang2020scorecam}.}
    \label{fig:scorecam_bridge}
\end{figure}


As shown in Fig.~\ref{fig:scorecam_bridge}, this image is an excellent example of enhanced structural integrity. With spatial awareness enabled, the model's attention is clearly and continuously distributed along the entire lateral structure of the bridge, indicating that the model is identifying the bridge as a complete, continuous object. With spatial awareness disabled, the attention becomes scattered, focusing only on one end of the bridge and some irrelevant background areas. This comparison strongly demonstrates that the spatially-aware module guides the model from focusing on scattered feature points to grasping the complete structure of the bridge.

\begin{figure}[!t]
    \centering
    \includegraphics[width=\linewidth]{Bridge_GradCAM_blocks10.png}
    \caption{Bridge class visualization using Grad-CAM (target layer: last ViT block attention; gradient-based localization; bilinear upsampling) with and without spatial awareness \cite{selvaraju2017gradcam}.}
    \label{fig:gradcam_bridge_new}
\end{figure}

Fig.~\ref{fig:gradcam_bridge_new} further demonstrates the comprehensive perception of structure and context. With spatial awareness enabled, the model focuses its attention on the T-shaped traffic hub structure formed by the connection of the bridge and the road. It not only identifies the bridge itself but also pays attention to the connecting roads, understanding that this is a composite traffic structure. With spatial awareness disabled, the attention is scattered and fails to form a structural understanding of the entire traffic hub. This proves that the spatially-aware module not only enhances the structural perception of a single-object, like bridge, but also extends to a more complex network of spatial relationships composed of multiple related objects (bridge, road, water body).

\begin{table*}[!t]
\centering
\caption{Modular ablation study results across all three datasets. The ASP baseline uses a larger variational encoder, while our methods are all based on a parameter-efficient deterministic spatial encoder.}
\label{tab:ablation_comprehensive}
\begin{tabular}{@{}lcccccccccc@{}}
\toprule
\multirow{2}{*}{Configuration} & \multirow{2}{*}{SAPM} & \multirow{2}{*}{NCAC} & \multirow{2}{*}{DPEMA} & \multicolumn{3}{c}{Accuracy (\%)} & \multirow{2}{*}{\makecell{Total \\ Params(M)}} & \multirow{2}{*}{\makecell{Trainable \\ Params(M)}} & \multirow{2}{*}{Time(ms)} & \multirow{2}{*}{\makecell{Storage \\ (MB)}} \\
\cmidrule(lr){5-7}
& & & & NWPU & UCM & MSTAR & & & & \\
\midrule
ASP Baseline & - & - & - & 83.00 & 80.09 & 69.06 & 173.6 & 87.8 & 8.49 & 662.44 \\
SAPM only & \checkmark & - & - & 83.47 & 82.25 & 72.75 & 91.8 & 6.0 & 6.19 & 351.07 \\
SAPM + NCAC & \checkmark & \checkmark & - & 83.02 & 83.59 & 74.88 & 91.8 & 6.0 & 6.13 & 351.07 \\
SAPM + DPEMA (no NCAC) & \checkmark & - & \checkmark & 84.10 & 83.10 & 74.10 & 91.8 & 6.0 & 6.15 & 351.07 \\
NCAC + DPEMA (no SAPM) & - & \checkmark & \checkmark & 83.60 & 82.80 & 73.20 & 173.6 & 87.8 & 8.49 & 662.44 \\
\textbf{Full System (Ours)} & \checkmark & \checkmark & \checkmark & \textbf{84.52} & \textbf{84.41} & \textbf{76.70} & \textbf{91.8} & \textbf{6.0} & \textbf{6.13} & \textbf{351.07} \\
\bottomrule
\end{tabular}
\end{table*}

\textbf{Ablation Study and Efficiency Analysis.}
The ablation studies first reveal the efficiency gains from our improvements to the ASP architecture. By replacing ASP's variational encoder with our deterministic spatially-aware encoder, namely SAPM, trainable parameters are drastically reduced from 87.8M to 6.0M, a reduction of 93.20\%, total parameters are reduced by 47.10\%, and inference latency is lowered by 27.80\%, see Table~\ref{tab:ablation_comprehensive}. This replacement also brings average accuracy gains across the three datasets. Furthermore, the full system's total storage is significantly lower than ASP's, 351.07\,MB versus 662.44\,MB, further demonstrating its deployability on resource-constrained platforms.

In terms of performance, we analyze the contributions of individual components. First, the NCAC module is particularly effective under sparse support, delivering new-class accuracy gains of up to 14.85\,pp on the UCM dataset, with an average improvement of 8.42\,pp across its incremental sessions, see Table~\ref{tab:new_class_learning}.

Second, to isolate module synergy, we evaluate two dual-module combinations, see Table~\ref{tab:ablation_comprehensive}: SAPM + DPEMA, without NCAC, and NCAC + DPEMA, without SAPM. Specifically, SAPM + DPEMA validates that DPEMA can mitigate inter-session drift, thus preserving SAPM's strong representational gains while improving stability. NCAC + DPEMA shows that, under EMA protection, NCAC can yield stronger plasticity than DPEMA alone via informed prototype initialization.

A noteworthy observation on NWPU-RESISC45 is that the \emph{SAPM+NCAC} combination, 83.02\%, performs slightly lower than \emph{SAPM only}, 83.47\%. We believe this indicates diminishing returns when NCAC is applied on top of already-strong SAPM features without the stabilizing DPEMA update. In this scenario, lacking EMA support, NCAC's single-step context modulation may not persist, causing it to underperform the simple mean prototype of SAPM-only.

Most importantly, this phenomenon disappears entirely in the full system, SAPM + NCAC + DPEMA, once DPEMA is introduced. The full system achieves the best average accuracy on all datasets, for example 84.52\% on NWPU. This fully corroborates the critical role of DPEMA in balancing retention and plasticity and validates the superior synergy of all components working together.

\subsection{New Class Learning Analysis}
\begin{table*}[!t]
\centering
\caption{Cross-dataset new-class learning analysis (Top-1 Accuracy, \%). Bold indicates TerraSAP (Ours).}
\label{tab:new_class_learning}
\begin{tabular}{@{}l|ccc|ccc|ccc@{}}
\toprule
\multirow{2}{*}{Session} & \multicolumn{3}{c|}{NWPU-RESISC45 (New Class)} & \multicolumn{3}{c|}{UCM (New Class)} & \multicolumn{3}{c}{MSTAR (New Class)} \\
\cmidrule(lr){2-4} \cmidrule(lr){5-7} \cmidrule(lr){8-10}
& ASP & OURS & Improvement & ASP & OURS & Improvement & ASP & OURS & Improvement \\
\midrule
S2 & 57.01 & \textbf{55.92} & -1.09 & 47.37 & \textbf{62.22} & +14.85 & 39.05 & \textbf{42.57} & +3.52 \\
S3 & 54.58 & \textbf{57.17} & +2.59 & 58.60 & \textbf{63.86} & +5.26 & 43.17 & \textbf{39.78} & -3.39 \\
S4 & 55.75 & \textbf{62.99} & +7.24 & 55.56 & \textbf{60.70} & +5.14 & 46.18 & \textbf{53.97} & +7.79 \\
S5 & 59.60 & \textbf{61.77} & +2.17 & - & - & - & 45.27 & \textbf{60.31} & +15.04 \\
S6 & - & - & - & - & - & - & 45.36 & \textbf{52.92} & +7.56 \\
S7 & - & - & - & - & - & - & 43.81 & \textbf{43.62} & -0.19 \\
\midrule
\textbf{Average} & 56.74 & \textbf{59.46} & \textbf{+2.72} & 53.84 & \textbf{62.26} & \textbf{+8.42} & 43.81 & \textbf{48.86} & \textbf{+5.05} \\
\bottomrule
\end{tabular}
\end{table*}


 A cross-dataset summary (Table~\ref{tab:new_class_learning}) shows consistent improvements in new-class learning. Specifically, we observe +2.72\,pp on NWPU-RESISC45 (S2--S5), +8.42\,pp on UCM, and +5.05\,pp on MSTAR (S2--S7). Occasional smaller gains or drops, for example, NWPU S2 and MSTAR S3/S7, typically arise when the estimated spatial context is noisy or atypical, in which case the prototype modulation strength may over- or under-correct; and when the DPEMA fusion occasionally favors stability late in the sequence. As mitigation, we clip the effective modulation via $\gamma_c \leftarrow \min(\gamma_c,\, \gamma_{\max})$ and adopt a robust aggregation for $\mathrm{ctx}_c$ to suppress outliers.


\subsection{Parameter Sensitivity Analysis}
To investigate the impact of key hyperparameters on model performance, we conducted a comprehensive one-factor-at-a-time, OFAT, sensitivity analysis across the NWPU-RESISC45, UCM, and MSTAR datasets. We focused on five core hyperparameters that govern spatial feature integration and the stability-plasticity balance. The results, summarized in Table~\ref{tab:sensitivity_analysis_table} and visualized in Fig.~\ref{fig:param_sensitivity_curves}, reveal consistent trends and provide clear guidance for optimal configuration.

\begin{table*}[!t]
\centering
\caption{Parameter sensitivity analysis of key hyperparameters on the three datasets. The best-performing value for each parameter is highlighted in bold.}
\label{tab:sensitivity_analysis_table}
{\setlength{\tabcolsep}{2.5pt}\scriptsize
\begin{tabular}{@{}l c ccc c ccc c ccc@{}}
\toprule
\multirow{2}{*}{\textbf{Parameter}} & \multicolumn{4}{c}{\textbf{NWPU-RESISC45}} & \multicolumn{4}{c}{\textbf{UCM}} & \multicolumn{4}{c}{\textbf{MSTAR}} \\
\cmidrule(lr){2-5} \cmidrule(lr){6-9} \cmidrule(lr){10-13}
& \textbf{Value} & \textbf{Avg Acc. (\%)} & \textbf{New (\%)} & \textbf{Old (\%)} & \textbf{Value} & \textbf{Avg Acc. (\%)} & \textbf{New (\%)} & \textbf{Old (\%)} & \textbf{Value} & \textbf{Avg Acc. (\%)} & \textbf{New (\%)} & \textbf{Old (\%)} \\
\midrule
\multirow{3}{*}{Modulation Strength ($\eta$)}
& 0.00 & 83.91 & 61.15 & 86.88 & 0.00 & 83.85 & 61.50 & 86.90 & 0.00 & 73.68 & 48.81 & 77.50 \\
& \textbf{0.10} & \textbf{84.52} & \textbf{61.77} & \textbf{91.76} & \textbf{0.10} & \textbf{84.41} & \textbf{62.26} & \textbf{87.60} & \textbf{0.10} & \textbf{76.70} & \textbf{43.62} & \textbf{82.86} \\
& 0.20 & 84.14 & 61.88 & 87.05 & 0.20 & 84.02 & 61.90 & 87.15 & 0.20 & 74.23 & 49.15 & 78.01 \\
\midrule
\multirow{3}{*}{Spatial Relation Weight  ($\omega_{\text{rel}}$)}
& 0.20 & 83.79 & 60.91 & 86.80 & 0.20 & 83.90 & 61.60 & 86.95 & 0.20 & 71.28 & 47.50 & 75.10 \\
& \textbf{0.35} & \textbf{84.52} & \textbf{61.77} & \textbf{91.76} & \textbf{0.35} & \textbf{84.41} & \textbf{62.26} & \textbf{87.60} & \textbf{0.35} & \textbf{76.70} & \textbf{43.62} & \textbf{82.86} \\
& 0.50 & 84.11 & 61.85 & 87.03 & 0.50 & 84.15 & 61.95 & 87.25 & 0.50 & 71.37 & 47.65 & 75.25 \\
\midrule
\multirow{5}{*}{Avg Alpha ($\alpha_{\text{avg}}$)}
& 0.30 & 83.95 & 61.10 & 86.95 & 0.30 & 83.50 & 61.10 & 86.50 & 0.30 & 72.02 & 47.90 & 75.90 \\
& 0.60 & 84.11 & 61.50 & 87.10 & 0.60 & 83.90 & 61.80 & 86.90 & 0.60 & 71.06 & 47.10 & 75.00 \\
& 0.75 & 83.99 & 62.10 & 86.85 & \textbf{0.75} & \textbf{84.41} & \textbf{62.26} & \textbf{87.60} & \textbf{0.75} & \textbf{76.70} & \textbf{43.62} & \textbf{82.86} \\
& 0.85 & 83.83 & 61.90 & 86.70 & 0.85 & 83.80 & 61.70 & 86.80 & 0.85 & 61.54 & 41.50 & 65.50 \\
& \textbf{0.95} & \textbf{84.52} & \textbf{61.77} & \textbf{91.76} & 0.95 & 83.95 & 61.90 & 87.00 & 0.95 & 71.39 & 47.50 & 75.30 \\
\midrule
\multirow{4}{*}{Dual EMA Fusion Weight ($\alpha_{\mathrm{ema}}$)}
& \textbf{0.35} & \textbf{84.52} & \textbf{61.77} & \textbf{91.76} & 0.35 & 84.10 & 61.95 & 87.15 & 0.35 & 69.63 & 46.50 & 73.50 \\
& 0.45 & 83.92 & 61.80 & 86.90 & \textbf{0.45} & \textbf{84.41} & \textbf{62.26} & \textbf{87.60} & \textbf{0.45} & \textbf{76.70} & \textbf{43.62} & \textbf{82.86} \\
& 0.55 & 84.10 & 62.00 & 87.05 & 0.55 & 84.00 & 61.90 & 87.00 & 0.55 & 67.76 & 45.50 & 72.00 \\
& 0.65 & 83.76 & 61.50 & 86.75 & 0.65 & 83.85 & 61.65 & 86.85 & 0.65 & 70.41 & 47.00 & 74.50 \\
\midrule
\multirow{5}{*}{EMA Beta ($\beta_{\text{ema}}$)}
& 0.980 & 84.10 & 61.95 & 87.05 & 0.980 & 83.95 & 61.80 & 87.00 & \textbf{0.980} & \textbf{76.70} & \textbf{43.62} & \textbf{82.86} \\
& \textbf{0.990} & \textbf{84.52} & \textbf{61.77} & \textbf{91.76} & \textbf{0.990} & \textbf{84.41} & \textbf{62.26} & \textbf{87.60} & 0.990 & 71.19 & 47.30 & 75.10 \\
& 0.995 & 84.03 & 61.90 & 87.00 & 0.995 & 84.00 & 61.85 & 87.05 & 0.995 & 67.95 & 45.80 & 72.20 \\
& 0.997 & 83.65 & 61.40 & 86.60 & 0.997 & 83.70 & 61.50 & 86.70 & 0.997 & 71.18 & 47.25 & 75.05 \\
\bottomrule
\end{tabular}}
\end{table*}


\begin{figure*}[!t]
    \centering
    \subfloat[Avg Alpha ($\alpha_{\text{avg}}$)]{\includegraphics[width=0.32\textwidth]{sensitivity_avg_alpha.png}}\hfill
    \subfloat[Dual EMA fusion weight ($\alpha_{\mathrm{ema}}$)]{\includegraphics[width=0.32\textwidth]{sensitivity_dual_ema_fusion_weight.png}}\hfill
    \subfloat[EMA $\beta_{\text{ema}}$]{\includegraphics[width=0.32\textwidth]{sensitivity_EMA_beta.png}}\\[1ex]
    \subfloat[Modulation strength ($\eta$)]{\includegraphics[width=0.32\textwidth]{sensitivity_modulation_strength.png}}\hfill
    \subfloat[Spatial relation weight ($\omega_{\text{rel}}$)]{\includegraphics[width=0.32\textwidth]{sensitivity_spatial_relation_weight.png}}
    \caption{Parameter sensitivity curves: x-axis shows parameter values and y-axis shows Avg Top-1 (\%). Each curve corresponds to a dataset (NWPU, UCM, MSTAR) with fixed color/marker across subfigures.}
    \label{fig:param_sensitivity_curves}
\end{figure*}

\subsubsection{\textbf{Modulation Strength ($\eta$)}} This parameter controls the intensity of the three-branch modulation applied to the spatial prompts. Across all datasets, performance peaks at a moderate value of $\eta=0.10$. A value of 0.00, which disables the modulation, results in a noticeable performance drop, indicating that the smoothing and stabilization from the modulation are beneficial. Conversely, increasing $\eta$ to 0.20 also leads to a slight decline, suggesting that excessive modulation may blur essential spatial details.

\subsubsection{\textbf{Spatial Relation Weight($\omega_{\text{rel}}$)}} This weight balances the contribution of the learned relative position bias in the self-attention mechanism of the SAPM. On NWPU, the peak is slightly higher around 0.50, while UCM and MSTAR favor 0.35; overall performance is stable within 0.35--0.50. Lower values, for example, 0.20, slightly degrade performance, implying that down-weighting the explicit spatial relationships is suboptimal. Higher values, for example, 0.50, show marginal or no improvement, suggesting that beyond a certain point, further emphasizing the spatial bias does not yield additional benefits and may interfere with content-based attention.

\subsubsection{\textbf{TIP--TSP Mixing ($\alpha_{\mathrm{avg}}$)}} This coefficient mixes the historical average prompt (TIP) with the current task-specific prompt (TSP) before injection. The results show a convex trend, with optimal performance generally clustering around 0.75 to 0.95. Very low values hinder stable prompting from few shots, while very high values can over-smooth task-specific details.

\subsubsection{\textbf{Dual EMA Fusion Weight ($\alpha_{\mathrm{ema}}$)}} This parameter determines the balance between the stability path and the plasticity path in the DPEMA module. The optimal value varies slightly across datasets, with NWPU favoring a lower value (0.35) and MSTAR/UCM performing best at 0.45. This suggests that the ideal balance is dataset-dependent. Nonetheless, the performance is relatively stable within the [0.35, 0.55] range, indicating that a reasonably balanced fusion is key. Extreme values in either direction would likely compromise the stability-plasticity trade-off.

\subsubsection{\textbf{EMA Beta ($\beta_{\text{ema}}$)}} This momentum coefficient controls the update rate of the historical knowledge bases in DPEMA. Across datasets, the best performance appears with a high momentum value of $\beta_{\text{ema}}$ (see Table~\ref{tab:sensitivity_analysis_table}). For NWPU and UCM, the optimum is 0.990, while for MSTAR it is 0.980. Values that are too high, for example, $\geq$0.995, suppress plasticity, while lower values harm old-class retention.

In summary, our analysis shows that TerraSAP's performance is robust within reasonable ranges for its key hyperparameters. The consistent trends across diverse datasets confirm that the underlying mechanisms for spatial awareness and knowledge retention are effective and generalizable. Despite these strengths, our framework still exhibits several limitations, detailed below.

\subsection{\textcolor{blue}{Limitations and Failure Mode Analysis}}

\begin{table*}[t]
\centering
\caption{\textcolor{blue}{Detailed confusion matrix (\%) for NWPU S2 New Classes (C25--C29). The columns represent the top confounding Base Classes.}}
\label{tab:nwpu_s2_confusion}
{\setlength{\tabcolsep}{2.5pt}\scriptsize
\renewcommand{\arraystretch}{1.2}
\begin{tabular}{@{}lc ccccccccccccc@{}}
\toprule
\multirow{2}{*}{\textbf{True Label}} & \multirow{2}{*}{\textbf{Acc.}} & \multicolumn{12}{c}{\textbf{Prediction Distribution on Key Classes (\%)}} & \multirow{2}{*}{\textbf{Others}} \\
\cmidrule(lr){3-14}
 & & \textbf{Self} & \makecell{\textbf{C1}\\\textbf{Airport}} & \makecell{\textbf{C7}\\\textbf{Church}} & \makecell{\textbf{C10}\\\textbf{Comm. Area}} & \makecell{\textbf{C12}\\\textbf{Desert}} & \makecell{\textbf{C13}\\\textbf{Forest}} & \makecell{\textbf{C14}\\\textbf{Freeway}} & \makecell{\textbf{C16}\\\textbf{Ground Track}} & \makecell{\textbf{C18}\\\textbf{Indus. Area}} & \makecell{\textbf{C22}\\\textbf{Meadow}} & \makecell{\textbf{C24}\\\textbf{Mobile Home}} & \makecell{\textbf{C26}\\\textbf{Railway}} & \\
\midrule
\textbf{C25 Parking Lot} & \textbf{74.1} & \textbf{74.1} & 1.0 & 0 & 0 & \textbf{10.9} & 4.5 & 0 & 0 & 0 & 3.0 & 0 & 0 & 6.5 \\
\textbf{C26 Railway} & \textbf{72.6} & \textbf{72.6} & 1.0 & 0 & 1.0 & 0 & 0 & \textbf{18.9} & 0 & 0 & 0 & 0 & - & 6.5 \\
\textbf{C27 Railway Stn.} & \textbf{15.9} & \textbf{15.9} & \textbf{8.0} & \textbf{35.3} & \textbf{18.9} & 0 & 0 & 0 & \textbf{9.0} & 5.0 & 0 & 0 & 1.0 & 6.9 \\
\textbf{C28 Rec. Farmland} & \textbf{68.7} & \textbf{68.7} & 2.0 & 0 & 1.5 & 0 & 0 & 2.0 & 1.5 & 4.0 & 0 & 6.5 & 0 & \textbf{13.8} \\
\textbf{C29 River} & \textbf{48.8} & \textbf{48.8} & 1.5 & 1.5 & 2.5 & 0 & 0 & \textbf{37.3} & 0 & 3.5 & 0 & 0 & 3.5 & 1.4 \\
\bottomrule
\end{tabular}}
\end{table*}

\textcolor{blue}{Despite the overall superior performance of TerraSAP, we observed instances of small or negative performance gains in specific sessions, such as a $-1.09\%$ drop in new class accuracy on NWPU S2 and fluctuations in MSTAR S3/S7. To investigate the underlying causes, we conducted an in-depth failure-mode analysis correlating performance with image characteristics, supported by the confusion matrix for NWPU S2 new classes in Table~\ref{tab:nwpu_s2_confusion}. The analysis reveals two primary limitations of our spatially-aware approach:}

\subsubsection{\textcolor{blue}{Ambiguity in Complex Structural Semantics}}
\textcolor{blue}{The explicit injection of spatial priors can lead to confusion when new classes share highly similar topological layouts or building densities with existing base classes. As shown in Table~\ref{tab:nwpu_s2_confusion}, \textbf{Railway Station} achieves the lowest accuracy of 15.9\% and is most frequently misclassified as \textbf{Church} (35.3\%) or \textbf{Commercial Area} (18.9\%). This is likely because these categories share complex, large-scale building structures surrounded by paved open spaces. The spatial encoder appears to extract overlapping structural features (e.g., large central buildings), causing the model to struggle in distinguishing the specific semantic function of the facility.}

\subsubsection{\textcolor{blue}{Topological Prior Overfitting}}
\textcolor{blue}{In some cases, the model relies too heavily on geometric shapes (spatial priors) at the expense of texture. A striking example is the \textbf{River} class, which achieves 48.8\% accuracy but suffers from a significant 37.3\% confusion with \textbf{Freeway}. Both classes exhibit strong linear, winding spatial topologies. The spatial prompt likely emphasizes this "linear path" structure, causing the model to misidentify water bodies as roads, overriding the spectral textural differences. Similarly, \textbf{Railway} (72.6\%) shows confusion with \textbf{Freeway} (18.9\%) due to their shared linear characteristics, and \textbf{Parking Lot} (74.1\%) is occasionally confused with \textbf{Desert} (10.9\%) due to the similarity of large, flat, featureless surfaces.}

\textcolor{blue}{In contrast, classes with relatively distinct combinations of geometry and texture perform better. \textbf{Rectangular Farmland} achieves a robust 68.7\% accuracy, with minor confusion scattered across structurally similar classes like Mobile Home (6.5\%) or Industrial Area (4.0\%). This confirms that TerraSAP excels when spatial topology provides a unique signature but requires further refinement to balance spatial geometry with fine-grained textural discrimination.}

\section{CONCLUSION}

To address the unique challenges in remote sensing FSCIL from complex spatial relationships and large intra-class variation, we propose TerraSAP, a parameter-efficient synergistic framework. With SAPM, NCAC, and DPEMA, it attains a balance between accuracy and efficiency. Experiments show state-of-the-art or competitive results on multiple benchmarks, compared with the original ASP model, TerraSAP reduces trainable parameters by 93.20\% and improves inference speed by 27.80\%. We validate that integrating spatial priors into prompt learning improves recognition of sparse new classes and offers a practical solution for incremental remote sensing.

TerraSAP achieves the largest gains when classes have clear, multi-object spatial topology and informative context. Gains shrink for categories driven mainly by spectral or texture cues or for single-object SAR with little surrounding context; very low resolution or extremely scarce shots also limit performance.

Future work will strengthen transfer across modalities through lightweight, modality-specific tuning; develop adaptive strategies that downweight spatial cues when context is, for example, texture or frequency-aware prompting and learned gating; and explore higher resolution features or light backbone tuning for texture-centric classes. We also plan a session-aware approach to adjust the stability and plasticity trade-off, and to extend spatial prompting to detection and segmentation. We will release the code and configuration files upon acceptance.

% References (BibTeX + IEEEtran)
\bibliographystyle{IEEEtran}
\bibliography{refs}

\begin{IEEEbiography}[{\includegraphics[width=1in,height=1.25in,clip,keepaspectratio]{author_JiaXingZeng.jpg}}]{Jiaxing Zeng}
received the B.S. degree in Computer Science and Technology from Beijing Normal University, Zhuhai, China, in 2023. He is with the School of Computer and Electronic Information, Guangxi University, Nanning 530004, China, and is pursuing the M.S. degree in Intelligent Technology with the Faculty of Innovation Engineering, Macau University of Science and Technology, Macau, expected in 2026. His research interests include remote sensing image analysis and class-incremental learning.
\end{IEEEbiography}
\begin{IEEEbiography}[{\includegraphics[width=1in,height=1.25in,clip,keepaspectratio]{author_YifengTan.png}}]{Yifeng Tan}
(Student Member, IEEE) received a B.S. degree in Computer Science and Technology from the Guangxi University for Nationalities, Guangxi, China, in 2020, and is pursuing his Ph.D. degree in Computer Science and Technology at Guangxi University, expected in 2026. His research interests include hyperspectral image processing, hyperspectral image classification, pattern recognition.
\end{IEEEbiography}
\begin{IEEEbiography}[{\includegraphics[width=1in,height=1.25in,clip,keepaspectratio]{author_LinaYang.png}}]{Lina Yang}
(Member, IEEE) received the Ph.D. degree in software engineering from the University of Macau, Macau, in 2015. She is an associate professor at the School of Electrical Engineering of Guangxi University and the School of Computer and Electronic Information of Guangxi University, working in Nanning, China. Her research interests include hyperspectral image classification, hyperspectral image reconstruction, and image processing.
\end{IEEEbiography}
\begin{IEEEbiography}[{\includegraphics[width=1in,height=1.25in,clip,keepaspectratio]{author_SiweiZHang.jpg}}]{Siwei Zhang}
received the B.S. degree in Electronic Information Engineering from the Harbin Institute of Technology, Harbin, China, in 2009, the M.S. degree in Mechanical Engineering from the Guangdong University of Technology, Guangzhou, China, in 2015, and the Ph.D. degree in Computer Technology and Application from the Macau University of Science and Technology, Macau, in 2020. He is currently an Assistant Professor with the Institute of Systems Engineering and the Department of Engineering Science, Faculty of Innovation Engineering, Macau University of Science and Technology. His research interests include discrete event systems, Petri nets, scheduling, and control.
\end{IEEEbiography}
\begin{IEEEbiography}[{\includegraphics[width=1in,height=1.25in,clip,keepaspectratio]{author_LianhuiLiang.png}}]{Lianhui Liang}
(Member, IEEE) received his M.S. and Ph.D. degree from the College of Electronics and Information Engineering, Hunan University, Changsha, China, in 2019 and 2024. From 2022 to 2024, he was a visiting Ph.D. student at the Hyperspectral Computing Laboratory, University of Extremadura, Extremadura, Spain, supported by the China Scholarship Council. He is currently an assistant professor at the School of Electrical Engineering, Guangxi University. His research interests include hyperspectral image processing, signal processing, pattern recognition, and machine learning.
\end{IEEEbiography}
\vfill
\end{document}
